\GXNSFBodyText

\noindent\textbf{总体研究方案}\par

\noindent\textbf{技术路线}\par
项目依次开展公开资料核验、事件链构建、多模态表征、状态识别、扰动分析和跨案例验证。每条记录保留来源、获取时间、许可边界和人工校验状态，使模型结果能够回溯至相应观测证据。佐佐木希案例用于贯通方法链，同构人物对照用于检验完整模型，广西—东盟样本则用于考察相关方法组件的跨语言适用性。

\begin{table}[htbp]
  \centering
  \fontsize{10}{14}\selectfont
  \caption{研究模块、主要输入与可检验输出}
  \begin{tabularx}{\textwidth}{p{0.19\textwidth}X X}
    \toprule
    研究模块 & 主要输入 & 可检验输出 \\
    \midrule
    事件链构建 & 官方页面、作品信息、公开新闻稿 & 事件编码一致性、来源覆盖率、冲突记录 \\
    状态识别 & 文本、图像元数据、更新节奏、平台关系 & 变点位置、状态稳定性、替代模型差异 \\
    韧性估计 & 扰动窗口与状态序列 & 恢复时间、路径变化、指标敏感性 \\
    分层验证 & 同构人物对照与广西—东盟组件样本 & 模型外部效度、组件误差、失效边界 \\
    \bottomrule
  \end{tabularx}
  \label{tab:gx-modules}
\end{table}

\noindent\textbf{数据获取与质量控制}\par

\noindent\textbf{公开数据纳入规则}\par
数据源分为官方资料、具名媒体公开报道、作品数据库和平台公开元数据四级。涉及生日、出道年份、作品和频道开设等基础事实时，优先使用经纪公司或项目官方来源；涉及争议事件时，记录公开报道的时间、内容及本人是否回应，并保留原文摘要与链接，以降低二次转述造成的事实漂移。

\noindent\textbf{人工编码与一致性检验}\par
两名编码者独立标注媒介类型、角色类别、事件方向和来源等级。先在训练样本上统一规则，再在保留样本上计算 Cohen's $\kappa$；低一致性类别返回规则修订，不通过提高样本量掩盖概念含混。时间冲突采用“多源并列—置信度标记—暂不裁决”的策略处理。

\noindent\textbf{数据伦理与合规}\par
项目仅研究公开职业信息，不建立人脸识别库，不抓取非公开账号，不保存普通用户可识别信息，不对个人敏感属性进行预测。平台数据获取遵守服务条款、访问频率和版权要求；可发布数据集仅包含来源索引、派生特征和合成示例，原始受版权保护内容不随项目公开。

\noindent\textbf{模型构建与检验}\par

\noindent\textbf{多模态表征与平台校准}\par
文本侧使用多语言句向量与主题演化特征，视觉侧只使用获许可图片的低维风格特征或公开元数据，行为侧使用更新间隔、媒介切换和合作类型。模型设置平台专属校准映射和共享条件状态转移层，通过域对抗或分层校准减少平台口径差异，不对同一观测同时建立重复的生成式分解。

\noindent\textbf{状态识别与变点检测}\par
采用隐马尔可夫模型作为可解释基线，以贝叶斯变点检测识别结构突变，再与时序深度模型比较。评价指标包括留出时间窗预测误差、状态持续性、跨初始化一致性和专家盲评可解释性。对照样本的纳入规则、留出时间窗和模型失稳判据将在分析前固定，并报告不支持迁移的样本；若复杂模型未能稳定优于基线，则保留更简单模型。

\noindent\textbf{扰动关联与稳健性分析}\par
平台切换、职业节点等内生事件只做描述性时序比较。对具有预先界定时点的候选外生事件，先检查处理时点唯一性、对照可比性与前趋势，再将中断时间序列、合成控制或匹配对照作为稳健性比较；任何假设不满足时均只报告关联、路径差异和时间顺序。通过改变恢复窗口、指标权重、对照集合和缺失值机制检验结论稳定性，不以敏感性分析反向选择最有利权重。

\noindent\textbf{可行性分析}\par

\noindent\textbf{数据可行性}\par
佐佐木希的基本资料、跨媒介活动与官方频道节点均有公开来源，足以构造方法演示事件链。项目不以获得私人数据为前提；当平台历史数据缺失时，可使用公开存档、作品级元数据和合成扰动序列完成算法验证。

\noindent\textbf{方法可行性}\par
研究所需的事件抽取、多语言表征、时序分段和网络分析均有成熟基线。项目采用逐层验证策略：先保证规则编码可复核，再比较统计模型，最后评估复杂模型，避免技术堆叠。每个环节均设置简单基线和失败退出条件。

\noindent\textbf{风险及应对}\par
针对数据缺失与跨语言偏差，项目采用多源索引、缺失机制模拟、双语抽检和术语表进行控制；针对单案例外推和区域分析单位差异，分别设置同期同构人物对照和组件级探索性测试；针对高热度事件，则通过最小化采集、去标识化和结果分层降低伦理风险。
